%%%%%%%%%%%%%%%%%%%%%%%%%%%%%%
\section{Background}
\label{sec:motivation-background}

\paragraph{Multi-agent Execution.}
With the development of LLM applications, more and more tasks require the design of multi-agent collaboration to be completed. These multi-Agent applications divide user tasks into multiple subtasks and assign them to several Agents for parallel execution. Each Agent maintains its own conversation, advancing tasks through LLM requests and tool invocations: the model generates responses or initiates tool calls, with the returned results becoming input for subsequent requests. **Request** is a single model call, and an **Agent** can initiate multiple requests to complete subtasks, while **Workflow** connects these tasks through serial dependencies, parallel branches, and merge points. The application is responsible for workflow orchestration, and the inference service layer is responsible for executing model requests.
Figure~\ref{fig:agent-workflow-examples} illustrates this structure. In code generation, the main Agent assigns the CSV export function to the API, UI, test writing, and documentation Agents. The branches execute in parallel, alternating between model requests and file reading, editing, or checking. The main Agent then integrates the changes and runs tests. In deep research, the lead Agent assigns technical, cost, security, and supply investigations, and then combines the collected evidence into a report~\cite{anthropic2025multiagent}. The number of requests and tool calls required by each branch may vary. Therefore, after a model request is completed, its associated Agent and the entire workflow may still not be finished.

\begin{figure}[t]
  \centering
  \includegraphics[width=\columnwidth]{figures/agent-workflow-examples.pdf}
  \caption{Multi-Agent Workflow Examples.
    Each row is one agent: $R_i$ denotes its $i$th LLM request, with
    intervening tool results carried into the next request.
    Branches may have different request counts.
    The coordinator resumes after their results join.}
  \label{fig:agent-workflow-examples}
\end{figure}



\paragraph{Workflow completion time.}
Let $w$ represent a workflow, $a_w$ the time the workflow arrives at the system, and $c_w$ the time when all necessary outputs are ready. The total time from arrival to completion of the workflow is defined as the *Workflow Completion Time* (WCT), i.e., $\mathrm{WCT}_w=c_w-a_w$. It includes model execution, tool execution, queuing, and dependency waiting. In the code example, integration needs to wait for all necessary modifications. Delays on the last incomplete branch will delay integration, while branches that are completed earlier may absorb the same delay before merging. Request delays affect the WCT through these dependencies.




\paragraph{KV reuse in multi-agent workflows.}
The model requests on these branches repeatedly consume context from
earlier steps. Coding agents carry source files, edits and test outputs
into later requests; research agents accumulate retrieved documents,
intermediate findings and tool results. Processing this growing history
from scratch on every request would repeat work already performed.
During \emph{prefill}, an LLM computes per-layer keys and values (KV)
for input tokens; \emph{decoding} extends this state as output tokens
are generated. Retaining this KV allows a subsequent request with a
matching token prefix to compute only the new suffix, provided the model
and position semantics are compatible~\cite{gao2024cachedattention,zheng2024sglang}.

This reuse occurs both within and across agents. Successive requests
from one agent can reuse KV for their shared conversation history.
Different agents can also inherit a common prompt before adding their
own tasks. For example, the API, UI and test agents in
Figure~\ref{fig:agent-workflow-examples} can share KV for common project
instructions and repository context when these form the same initial
token sequence. Different role prompts or context edits can shorten
the reusable prefix~\cite{zheng2024sglang}.

Therefore, KV Cache is an important intermediate state in the process of Agent execution or LLM reasoning requests.


\begin{figure}[t]
  \centering
  \includegraphics[width=\columnwidth]{figures/gpu-server-availability.pdf}
  \caption{GPU-server availability reconstructed from InfiniteHBD's
    production-training trace~\cite{shou2025infinitehbd}: 400 servers minus
    those within recorded outage intervals. Records include GPU, network,
    power/cooling, and software faults; unknown causes, tests,
    troubleshooting, and changes are also retained. Overlapping intervals
    per server are merged. Availability is measured at the node level.}
  \label{fig:gpu-server-availability}
\end{figure}

\paragraph{Failures and KV reconstruction.}
These KV reuse benefits depend on the cached state remaining available.
In large-scale LLM deployments, hardware faults and software crashes can
interrupt serving workers and make their GPU-resident KV
unavailable~\cite{cao2026lumen,jayakody2026ghostserve}.
For production-cluster context, Figure~\ref{fig:gpu-server-availability}
shows the public InfiniteHBD training-cluster trace: 584 recorded outages
affect 231 of 400 GPU servers over approximately 349 days, with reconstructed
availability falling to 365 servers~\cite{shou2025infinitehbd}.
The trace includes faults and operational interruptions, as detailed in
the caption.

When no reusable KV survives, a replacement worker must reconstruct it
by replaying the retained tokens, even if it already has the model weights
and the application has preserved the conversation
history~\cite{strati2024dejavu}. For a long-running agent, this means
reprocessing the context accumulated across earlier requests rather than
only the latest input~\cite{jayakody2026ghostserve}.
One failed worker can also interrupt several co-located agents, requiring
surviving workers to reconstruct multiple histories while handling normal
traffic~\cite{cao2026lumen}. The resulting delays propagate through the
workflow's dependencies: reconstruction on the last unfinished branch
postpones its join and can increase WCT.

\paragraph{KV backup across storage tiers.}
To reduce reconstruction work, a \emph{KV checkpoint} retains an already
computed prefix. If it survives the failure, recovery retrieves this
prefix and replays only the unsaved
tail~\cite{strati2024dejavu,jayakody2026ghostserve}.
KV can be retained in GPU HBM, host DRAM, or shared
storage~\cite{gao2024cachedattention,qin2025mooncake}.
These tiers differ in both storage cost and recovery latency.
HBM offers fast access but has high capacity cost and competes with
foreground inference for GPU memory. DRAM and storage-backed shared
services can provide cheaper capacity, with additional retrieval latency.

Figure~\ref{fig:motivation-recompute} illustrates the latency difference
between reconstructing a history and reusing retained KV.
At 1,010,000 tokens, full prefill takes 571.21\,s even with sparse
attention; HBM, DRAM and shared-storage reuse take 1.305, 2.663 and
10.156\,s, respectively. The history occupies 57.92\,GB across two ranks.
All paths end when the same input's KV is ready: reuse restores the
preceding history, if needed, and computes one new token. It can be seen that through KV reuse, the delay of KV reconstruction can be greatly reduced. But we all know that storage has costs, and how to store and which key-value pairs to store has become an issue that cannot be ignored.

\begin{figure}[t]
  \centering
  \includegraphics[width=\columnwidth]{figures/recovery-cost.pdf}
  \caption{Measured KV reconstruction and reuse costs for
    Qwen2.5-7B-Instruct-1M on two A100 PCIe 40\,GB GPUs
    (TP2, BF16, batch one), with sparse attention above 32K.
    Full prefill starts from empty KV; reuse starts from $(C-1)$ cached
    tokens and computes one. KV reconstruction time includes any
    retrieval and computation until all $C$ tokens' KV is ready.}
  \label{fig:motivation-recompute}
  \Description{Four measured KV reconstruction-time curves for Qwen2.5-7B-Instruct-1M
    from 8K to 1,010,000 tokens: blue circles for resident HBM reuse,
    orange squares for DRAM reuse, purple diamonds for shared-storage
    reuse, and red triangles for full prefill. Both axes are logarithmic.
    At the longest context, the four paths take 1.305, 2.663, 10.156 and
    571.21 seconds, respectively.}
\end{figure}